\newpage
\topic{5 continued}{One shared face, one velocity correction}
\lessonfigure{shared_face_correction}{A shared-face pressure correction. When $p_E^\prime$ exceeds $p_P^\prime$ by 0.01, the velocity increment is westward: $u_e^\prime=-0.02$. The predictor 0.12 becomes 0.10. The same face velocity enters the two neighboring outward flux balances with opposite signs.}{shared-face-correction}
The predictor $u_e^*$ is obtained from momentum with the current pressure
estimate. It generally leaves a cell mass imbalance. A pressure correction
$p^\prime$ produces a velocity increment $u_e^\prime$:
\[
u_e=u_e^*+u_e^\prime,\qquad
u_e^\prime=-d_e\frac{p_E^\prime-p_P^\prime}{h},\qquad d_e=1/a_e>0.
\]
Here $a_e$ is the diagonal of the \emph{volume-normalized} momentum equation.
The prime denotes an increment, not differentiation. If $p_E^\prime>p_P^\prime$,
the correction acts westward, consistent with acceleration down a pressure
gradient.

Cell $P$ and cell $E$ use the same shared face value. In the outward flux
balance it contributes $+h u_e$ for $P$ and $-h u_e$ for $E$. A shared-face
correction therefore preserves pairwise cancellation between neighboring cells.
The pressure equation chooses these corrections to remove each cell's
predictor mass imbalance, subject to the boundary fluxes and pressure reference.

Return to the reading-map mass-balance example, where $u_e^*=0.12$ and
$r_{c,P}^*=0.08$. For this correction illustration, assume $a_e=2$, so
$d_e=0.5$, and suppose the pressure solve returns
$p_E^\prime-p_P^\prime=0.01$. Then
\[
u_e^\prime=-0.5(0.01)/0.25=-0.02,\qquad u_e=0.12-0.02=0.10.
\]
Holding the other faces fixed, cell $P$ loses $h(0.02)=0.005$ of outward
volume flux: its imbalance becomes zero. The same face change increases
cell $E$'s outward flux by $0.005$. A correction that fixes $P$ alone may
therefore spoil $E$; all pressure cells must be solved together.
The prescribed diagonal $a_e=2$ belongs to this correction illustration,
not to the earlier convection-flux example.
